% Lösungen Einheit 3 von 3 – Am Ratengraphen (zusammenbau v0.1)
\einheitenkopf{Einheit 3 von 3 · Am Ratengraphen}

\erg{14}{a) $r(1) = 3$: in diesem Moment fließen 3 Liter je Minute zu; Dreiecksfläche von 0 bis 4: Änderung $= 8$ Liter (Zunahme) \quad b) $∫_0^4 (6t - 1{,}5t^2)\,\mathrm{d}t = [3t^2 - 0{,}5t^3]_0^4 = 48 - 32 = 16$; nach vier Wochen ist die erste Pflanze 16\,cm höher als die zweite (nicht 16\,cm hoch) \quad c) $T = 6$: die Fläche über der Achse von 0 bis 3 ($4{,}5$) ist so groß wie die Fläche unter der Achse von 3 bis 6 ($4{,}5$) \quad d) Die Fläche unter e (I und II) ist die Zahl aller eingefahrenen Autos, die unter a (II und III) die aller ausgefahrenen. Weil das Parkhaus am Anfang und am Ende leer ist, sind beide gleich; nimmt man den gemeinsamen Teil II weg, bleibt I so groß wie III.}
\erg{15}{a) Tom nimmt gleiche Raten statt gleicher Bestände. Gleich viel heißt $∫_0^T (2t - 6)\,\mathrm{d}t = T^2 - 6T = 0$, also $T = 6$ \quad b) Die Fläche oberhalb ist die Zunahme, die unterhalb die Abnahme; der Bestand ist wieder gleich, wenn beide sich aufheben. \quad c) zum Beispiel die Gerade $r(t) = 3 - t$: Dreieck über der Achse von 0 bis 3 und gleich großes Dreieck darunter von 3 bis 6 \quad d) $200 + ∫_0^{10} (50 - 8t)\,\mathrm{d}t = 200 + 100 = 300\,\mathrm{m}^3$; am meisten bei $50 = 8t$, also $t = 6{,}25$}

15c)

\begin{ksys}[xmin=0,xmax=7,ymin=-4,ymax=4,klein] \funktionab{3-\x}{r}{0}{6} \end{ksys}

